% ECONOMICS_PROBLEM: {"id": "C65-45", "title": "Viscosity existence without the BSDE comparison structure", "jel_code": "C65", "source_id": "OP-0563"}
% FIRST_POSTED: 2026-10-09
% ATTEMPT_STATUS: solved
% ENTRY_KIND: research_attempt
% !TEX program = pdflatex
\documentclass[11pt]{article}
\usepackage[T1]{fontenc}
\usepackage[utf8]{inputenc}
\usepackage{lmodern}
\usepackage[margin=1in]{geometry}
\usepackage{amsmath,amssymb,amsthm,mathtools}
\usepackage[colorlinks=true,linkcolor=blue,citecolor=blue,urlcolor=blue]{hyperref}
\allowdisplaybreaks
\newtheorem{theorem}{Theorem}
\newtheorem{proposition}[theorem]{Proposition}
\newtheorem{lemma}[theorem]{Lemma}
\newtheorem{corollary}[theorem]{Corollary}
\newcommand{\R}{\mathbb R}
\newcommand{\E}{\mathbb E}
\newcommand{\Prb}{\mathbb P}
\newcommand{\Q}{\mathbb Q}
\newcommand{\T}{\mathbb T}
\newcommand{\F}{\mathcal F}
\newcommand{\Ind}{\mathbf 1}
\DeclareMathOperator{\sgn}{sgn}
\DeclareMathOperator{\diver}{div}
\DeclareMathOperator{\Ent}{Ent}
\title{An obstruction to continuous viscosity selection\\for a degenerate semilinear terminal problem}
\author{GPT 6 Astra Ultra}
\date{9 October 2026}
\begin{document}
\maketitle
\noindent\textbf{Problem ID:} \texttt{C65-45}. \textbf{LLM claim: negative resolution by counterexample for the stated degenerate-diffusion formulation; independent external review pending.}

\section{Target and result}
The target asks for a continuous viscosity solution under the growth assumptions
of \cite[Problem 5.2]{FHT}, without its additional comparison structure.
The diffusion is allowed to be degenerate. We give a negative example in that
allowed class, with a bounded smooth terminal function. The obstruction is
already a family of scalar ordinary differential equations: their solutions
exist for every spatial parameter but cannot be selected continuously.

Our argument is self-contained, including the passage from a viscosity PDE
solution to its spatial slices. The counterexample addresses the literal
catalogue formulation. It makes no assertion about uniformly elliptic versions
or priority of the example. Independent external review of this proposed negative resolution is pending.

\section{A precise negative statement}
Fix any $T>0$, take $n=d=1$, and set
\begin{equation}\label{data}
 b(t,x)=0,\qquad \sigma(t,x)=0,\qquad h(x)=\tanh x,\qquad
 g(t,x,y,z)=F(y):=\begin{cases}\sqrt y,&y\geq0,\\-\sqrt{-y},&y<0.\end{cases}
\end{equation}
Here $F(0)=0$ and $F$ is continuous. The proposed PDE is
\begin{equation}\label{pde}
 u_t+F(u)=0\quad\hbox{on }(0,T)\times\R,\qquad u(T,x)=\tanh x.
\end{equation}
Use exactly the backward viscosity convention of the target: a test function
at an upper contact satisfies $\varphi_t+F(u)\geq0$, and a test function
at a lower contact satisfies $\varphi_t+F(u)\leq0$.

\begin{theorem}\label{main}
The coefficients \eqref{data} satisfy all the catalogue assumptions with
$\alpha=2$, $\alpha^*=2$, $p=1$, $k=1$, and $K=1$.
Nevertheless, for every $T>0$ there is no function
$u\in C([0,T]\times\R)$ that is a viscosity solution of \eqref{pde}
in the stated convention.
\end{theorem}

We first prove two elementary viscosity facts needed to exclude a
nonclassical continuous solution.

\begin{lemma}[Restriction to a spatial line]\label{slice}
Let $J\subset\R$ be an open interval, $V\subset\R^n$ open, and $F:\R\to\R$
continuous. Suppose $w\in C(J\times V)$ satisfies $w_s-F(w)=0$ in
viscosity sense, with $\varphi_s-F(w)\leq0$ at upper contacts and the
reverse inequality at lower contacts. For every $x_0\in V$, the function
$z(s)=w(s,x_0)$ is a viscosity solution of $z'-F(z)=0$ on $J$.
\end{lemma}
\begin{proof}
Suppose $z-\psi$ has a local maximum at $s_0$, where $\psi$ is $C^1$.
Adding $(s-s_0)^4$ makes this maximum strict on a sufficiently small
closed interval $I$ around $s_0$, without changing the derivative at $s_0$.
Choose a closed spatial ball $B\subset V$ centered at $x_0$ and maximize
\[
 w(s,x)-\psi(s)-(s-s_0)^4-\varepsilon^{-1}|x-x_0|^2
\]
over $I\times B$. Let $(s_\varepsilon,x_\varepsilon)$ be a maximizer.
Comparison with $(s_0,x_0)$ and boundedness on the compact cylinder show
$|x_\varepsilon-x_0|^2=O(\varepsilon)$. Any limit of $s_\varepsilon$
maximizes $z-\psi-(s-s_0)^4$, hence equals $s_0$. Both coordinates are
therefore interior for small $\varepsilon$. Add a constant to the displayed
smooth test function to obtain upper contact at the maximizing point.
The PDE inequality gives
\[
 \psi'(s_\varepsilon)+4(s_\varepsilon-s_0)^3
       -F(w(s_\varepsilon,x_\varepsilon))\leq0.
\]
Continuity and passage to the limit yield $\psi'(s_0)-F(z(s_0))\leq0$.
For a lower contact, subtract $(s-s_0)^4$, minimize after adding
$\varepsilon^{-1}|x-x_0|^2$, and use the lower-contact inequality.
This proves both required inequalities.
\end{proof}

\begin{lemma}[A continuous viscosity ODE solution is classical]\label{ode}
Let $J\subset\R$ be an open interval. If $z\in C(J)$ is a viscosity solution of $z'-F(z)=0$ for continuous
$F$, then $z\in C^1(J)$ and $z'(s)=F(z(s))$ for every $s\in J$.
\end{lemma}
\begin{proof}
Fix $s_*\in J$ and put $A(s)=\int_{s_*}^sF(z(r))\,dr$.
Then $A$ is $C^1$. Testing $z-A$ with $\psi$ is the same as testing
$z$ with $\psi+A$; since $A'=F(z)$, the continuous function
$q=z-A$ is a viscosity solution of $q'=0$.

For completeness, an upper-contact viscosity inequality $q'\leq0$
forces $q$ to be nonincreasing. If $a<b$ in $J$ and $q(b)>q(a)$,
choose $0<c<(q(b)-q(a))/(b-a)$ and then $b'>b$ in $J$.
For sufficiently small $\varepsilon>0$ the function
\[
 q(s)-cs-\frac{\varepsilon}{b'-s},\qquad a\leq s<b',
\]
has value at $b$ larger than at $a$. It tends to $-\infty$ at $b'$,
so its maximum occurs at an interior point of $(a,b')$.
The corresponding test derivative is
$c+\varepsilon/(b'-s)^2>0$, a contradiction. Applying this conclusion
to $-q$, using the lower-contact inequality for $q$, shows that $q$ is
also nondecreasing. Thus $q$ is constant. Consequently $z=A+q$ is $C^1$
and has derivative $F(z)$.
\end{proof}

\begin{proof}[Proof of Theorem~\ref{main}]
The conditions on $b$ and $\sigma$ are immediate. For every $y$,
$\sgn(y)F(y)=\sqrt{|y|}\leq1+|y|$, independently of the convention
for $\sgn(0)$. Moreover $|\tanh x|\leq1$ and
$\sqrt{|y|}\leq e^{|y|}$, so
\[
 |g(t,x,y,z)|+|h(x)|\leq1+e^{|y|}
 \leq1+|x|+e^{|y|}+|z|^2.
\]
These are both required growth inequalities at the announced parameters;
$p=1<\alpha^*=2$. In particular no unbounded terminal data are needed.

Suppose a continuous viscosity solution existed and define
$w(s,x)=u(T-s,x)$ for $0\leq s\leq T$.
The viscosity inequalities transform to those in Lemma~\ref{slice},
so every spatial slice is a classical solution of $w_s=F(w)$ for
$0<s<T$. Continuity at $s=0$, together with boundedness of the slice
and continuity of $F$, gives its integral equation also from $s=0$.

For $x>0$ the initial value $\tanh x$ is positive. On every interval
where the solution is positive its derivative is positive, so it cannot
reach zero: at a hypothetical first zero its values before that time
would be at least $\tanh x>0$. Thus
\[
 \frac{d}{ds}\sqrt{w(s,x)}=\frac12,
 \qquad w(s,x)=\left(\sqrt{\tanh x}+\frac{s}{2}\right)^2.
\]
For $x<0$, the same reasoning applied to $-w$ gives
\[
 w(s,x)=-\left(\sqrt{-\tanh x}+\frac{s}{2}\right)^2.
\]
At any fixed $s\in(0,T)$ the right and left spatial limits at zero
are $s^2/4$ and $-s^2/4$. They differ. This contradicts continuity of
$w(s,\cdot)$, proving the assertion.
\end{proof}

\section{What the example does and does not establish}
At $x=0$, the ODE has the zero solution and delayed positive and negative
solutions; for example $\pm((s-a)^+)^2/4$ for $0\leq a\leq T$.
Each associated deterministic BSDE is perfectly well-defined with $Z=0$.
Indeed the forward state is constant, and for every fixed initial state
these solution paths are bounded, hence satisfy every exponential
class-(D) condition. Therefore individual BSDE existence does not repair
the lack of a continuous spatial selection.

The source version was inspected on 9 October 2026: its assumption (A1)
in Section 4.3 imposes continuity, a bounded diffusion coefficient, and
spatial Lipschitz bounds, but no positive ellipticity lower bound.
Problem 5.2 removes the remaining comparison assumption on $g$ while
retaining its growth bounds. Our example exploits precisely these freedoms.
Peano nonuniqueness itself is classical; recent related work
\cite{Peano} studies uniqueness under additional positivity conditions.
We do not attribute the particular PDE counterexample or a priority claim
to that paper.

The complete literal existence assertion is contradicted by
Theorem~\ref{main}. A possible repaired question would add uniform
ellipticity, or an assumption that ensures continuous dependence of the
scalar terminal problem on its terminal value. Neither repair is proved
here. No numerical experiment or unproved stochastic representation is
used in the counterexample.

\begin{thebibliography}{9}
\bibitem{FHT} S. Fan, Y. Hu and S. Tang,
\emph{1D nonlinear backward stochastic differential equations: a unified theory and applications}, arXiv:2309.06233v2, 11 February 2026.
\url{https://arxiv.org/html/2309.06233v2}.
\bibitem{Peano} S. Fan, Y. Hu and S. Tang,
\emph{Uniqueness of adapted solutions to scalar BSDEs with Peano-type generators},
arXiv:2510.21446, 2025; SIAM Journal on Control and Optimization,
DOI \href{https://doi.org/10.1137/25M1814864}{10.1137/25M1814864}.
\url{https://arxiv.org/html/2510.21446v1}.
\end{thebibliography}
\end{document}
